\documentclass[11pt]{article}

\usepackage[utf8]{inputenc}
\usepackage[T1]{fontenc}
\usepackage{lmodern}
\usepackage[margin=1in]{geometry}
\usepackage{microtype}
\usepackage{setspace}
\usepackage{natbib}
\usepackage{amsmath,amssymb}
\usepackage{enumitem}
\usepackage[hidelinks]{hyperref}
\usepackage{titlesec}

\setstretch{1.08}
\setlength{\parindent}{1.35em}
\setlength{\parskip}{0pt}
\setlist[itemize]{leftmargin=1.7em}
\setlist[enumerate]{leftmargin=1.7em}
\titleformat{\section}{\normalfont\large\bfseries}{\thesection.}{0.5em}{}
\titleformat{\subsection}{\normalfont\normalsize\bfseries}{\thesubsection.}{0.5em}{}

\title{The Incoherence of ``Possibility'':\\ Why Machine Consciousness Remains Barred in Principle}
\author{Anonymous}
\date{2026}

\begin{document}
\maketitle

\begin{abstract}
\noindent The recent defense of machine consciousness's possibility, while admirably restrained, rests on a foundational confusion between the removal of certain objections and the establishment of genuine modal possibility. I argue that the paper's central strategy---the defeater-filter test---not only fails to establish what it claims but actively presupposes the very assumption it is meant to defend: namely, that consciousness can be multiply realized. The paper's two main positive arguments, the brain-in-a-vat case and the organizational-invariance thesis, both collapse under scrutiny. The vat case is not consciousness-preserving in the relevant sense; it preserves biological continuity while removing the embedding that may constitute the ground of consciousness. Organizational invariance, even granting its premises, applies only to functional duplicates embedded in the right biological context---a conclusion that supports biological essentialism, not machine consciousness. The paper's concluding modesty masks a failure to engage with the deepest level of the debate: whether consciousness is intrinsically tied to the kind of \textit{integration} characteristic of living systems, not merely to formal organization. I contend that the answer is yes, and that no amount of functional organization, divorced from biological development, embodied coupling, and metabolic integration, will suffice for consciousness. Machine consciousness is not merely unproven; it is in principle impossible given what we understand about the nature of conscious life.
\end{abstract}

\section{The Confusions Underlying ``Possibility''}

The paper's opening move---distinguishing four senses of ``possible'' and claiming to establish only epistemic possibility---is presented as methodological virtue \citep{intermittent}. In fact, it masks a fatal ambiguity that runs through the entire argument.

Epistemic possibility, as the author defines it, is that ``we do not know that machine consciousness is impossible.'' This is trivially true about almost everything. We do not know that rocks are not conscious, or that the color red does not have phenomenal experiences, or that mathematical operations have inner lives. Epistemic possibility is so weak that establishing it tells us virtually nothing. The paper's claimed achievement---that nothing demonstrated rules machine consciousness out---is true but uninteresting without a positive argument for expecting it.

The author promises that ``the upgrade to nomological or metaphysical possibility rests on the positive argument of Section 6.3, the organizational-invariance argument.'' But this is precisely where the argument fails. Organizational invariance, even if granted in full, establishes at most that \textit{some} substrate is not necessary for \textit{some} systems. It does not establish that consciousness could arise in a purely computational substrate divorced from biological embedding, metabolic integration, developmental history, and the kind of embodied coupling with an environment that characterizes living systems.

The central confusion is this: the paper treats ``we cannot rule it out'' as equivalent to ``it is possible.'' But these are fundamentally different claims. That we cannot provide a deductive proof of impossibility does not entail that we have any reason to expect possibility, nor does it mean the burden of proof has shifted to the skeptic. When facing an unprecedented claim---machine consciousness---the default epistemological stance should be: absence of evidence that $X$ can occur without feature $Y$ is evidence of absence unless we have positive theoretical reason to expect $X$ in the absence of $Y$. The paper provides no such positive theoretical reason; it merely removes some alleged barriers while presupposing the removal is complete.

\section{The Defeater-Filter Test and Its Presupposition}

The paper's diagnostic test---``would this objection unseat a human under consciousness-preserving conditions?''---is its methodological centerpiece. But the test presupposes precisely what needs to be established.

Consider how the test works. An objector says: ``Machines have feature $X$ (intermittency, discreteness, resettability) which disqualifies them from consciousness.'' The filter responds: ``But humans could have feature $X$ under consciousness-preserving conditions, so $X$ is not intrinsically disqualifying.'' This conclusion follows \textit{if and only if} we grant that human consciousness would be preserved under those conditions. But what justifies that grant?

For the vat-brain case, the paper argues that intermittency is harmless because we accept it in sleep. This works because we have independent, overwhelming evidence that consciousness is preserved through sleep. We can sleep and wake with full continuity of memory, identity, and body. We know this from first-person experience. But for machines, we have no such evidence. We have no data whatsoever about whether consciousness would be preserved if a machine were intermittent, because we have no baseline evidence that machines are conscious in the first place.

The defeater-filter test, therefore, illicitly trades on the asymmetry in our evidence. For humans, we use the test to show that certain features are not \textit{disqualifying per se} because we already know humans are conscious despite those features. For machines, we attempt to use the test to show that the same features are not disqualifying---but we have no baseline knowledge that machines are conscious at all. The test cannot do the work the paper claims because it requires an evidential asymmetry it tries to dissolve.

More precisely: the test is sound \textit{within a context where we are confident about consciousness}, but it becomes question-begging \textit{when deployed across a boundary of radical uncertainty}. The paper acknowledges that the test ``is not a positive criterion of consciousness'' and that it ``is silent on substrate, embodiment, and intentionality.'' Exactly so. And it is precisely on substrate, embodiment, and intentionality that the case for machine consciousness fails.

The paper's own language reveals the problem. It states that architectural features are ``orthogonal to consciousness \textit{relative to a consciousness-capable system}.'' This is a transparent retreat. The question is not whether features are orthogonal to consciousness once we have a consciousness-capable system; the question is what \textit{makes} a system consciousness-capable in the first place. By restricting the test's application to systems we already know are consciousness-capable, the paper ensures the test cannot arbitrate the very question at issue.

\section{The Vat-Brain Case: Preserving Biology, Not Consciousness}

Section 2's brain-in-a-vat argument is intriguing but ultimately question-begging in a more subtle way. The case stipulates a human \textit{brain}, not a human \textit{person} or a human \textit{organism}. This is not an innocent choice; it is where the begging occurs.

A human brain suspended in a vat and intermittently connected to a synthesizer is certainly still a human brain. If human brains are intrinsically conscious substrates---if, as biological naturalism suggests, consciousness is a causal power of specifically evolved biological tissue \citep{searle1992}---then the vat-brain remains conscious by virtue of \textit{remaining biological tissue}. But this preservation of consciousness proves nothing about machines, because machines are not biological tissue. The case merely shows that \textit{if consciousness is a capacity of the tissue itself}, then the tissue remains conscious despite architectural manipulation. This is entirely consistent with biological essentialism.

The paper acknowledges, correctly, that the vat case ``holds the contested variable fixed'' by stipulating a human brain. And it states: ``The one variable the vat does not neutralize---because it holds it fixed by stipulating a human brain---is substrate, to which the second half of the paper is owed.'' This is honest. But it means the vat argument does not address substrate at all. It merely shows that intermittency, reset, and reboot do not \textit{automatically} unseat consciousness in a biological substrate. From this, no inference to machines follows.

Furthermore, the paper's treatment of reset is incomplete. The author writes: ``loss of accumulation is not loss of momentary experience,'' drawing on the Clive Wearing case. But Wearing is not an example of consciousness despite loss of accumulation; he is an example of consciousness \textit{in a living person} despite impaired \textit{consolidation}. The crucial disanalogy is this: Wearing's moment-to-moment consciousness occurs within a unified, metabolically integrated, developing organism whose neural tissue has the continuous causal history of a living brain. A snapshot-reset system lacks this history. It is not merely a system with memory loss; it is a system that is \textit{recreated} from a template.

The paper notes this disanalogy carefully, distinguishing ``explicit episodic memory, implicit memory, affective continuity, embodied skill, and total state-restoration.'' But after noting these distinctions, it does not adequately explore their relevance. Affective continuity may depend on embodied skill, which may depend on a continuous metabolic substrate. Implicit memory itself may be inseparable from the gradual sculpting of neural tissue over time. None of these are straightforwardly preserved in a reset system.

More fundamentally: the vat case preserves something the paper does not name carefully enough---the \textit{developmental and evolutionary history of the biological tissue}. A brain suspended in a vat may lack input/output contact with its original body, but it retains the structure and chemical gradients developed over millions of years of evolution and shaped through an individual's personal history. A digital system reset to a prior snapshot retains no such history. It is, at best, a replica, not a continuation. And if consciousness depends on the developmental and historical integration of a biological system, then the vat case proves nothing about machines, because machines lack that history.

\section{Discrete Time and the Disguised Assumption}

Section 3's argument about discrete time is elegant but contains a concealed sleight of hand.

The paper grants the conditional: \textit{if} time is discrete, then consciousness is compatible with discreteness. From this, it argues that subjective continuity can float free of objective continuity, and that gaps are ``experientially invisible'' because nothing experiences a gap during which it is not running. This is sound reasoning \textit{within the hypothetical}. But the paper then moves from this conditional to a claim about three cases, including digital machines, as if the result transfers.

The argument's critical move is stated carefully: ``three cases argue well for (a) [physical discreteness does not preclude consciousness] and (b) [subjective continuity need not mirror objective continuity], suggestively for (c) [invisible pauses do not necessarily interrupt experience from the inside]---and not at all for (d) [digital computation satisfies the same consciousness-relevant temporal conditions as a brain].'' The paper explicitly disclaims establishing (d).

But here is the problem: the entire force of the section depends on cases (i) and (iii)---both of which involve \textit{human brains}. Case (ii), the RL system, is presented as a parallel candidate, but the argument's persuasiveness rests entirely on the reader accepting that human brains could be discrete without losing consciousness. This carries no implications for machines.

Furthermore, the argument obscures a crucial distinction: between a \textit{physical substrate} being discrete and a \textit{computational process} being discrete. If time is fundamentally discrete, then everything in the universe, including rocks and plants and electrons, unfolds in discrete ticks. Discreteness in this sense is a background condition, not a feature that distinguishes machines from organisms. But \textit{digital computation} is not merely discrete in this sense; it is discrete \textit{at the level of formal state transitions}. A discrete physical process---a biological brain evolving tick by tick in a discrete time---maintains causal continuity at the physical level. Its states at tick $n$ stand in appropriate physical causal relations to states at tick $n+1$, mediated by the continuous physical gradients and interactions of neural tissue.

A digital computer, by contrast, may have states at tick $n$ that stand in no direct causal relation to states at tick $n+1$ except through \textit{representation}---the interpretation of physical substrate \textit{as} computing states. This is the distinction the paper correctly identifies but fails to adequately emphasize: ``physical discreteness is not digital computation.'' If time itself is discrete, that is bad news for machines, not good news. It means consciousness in us despite discrete time proves nothing about whether consciousness arises in systems whose organization is constitutively digital-computational rather than merely substrate-discrete.

The deepest issue is this: if time is discrete, the physical causal relationships between our neural states at adjacent ticks are still relationships between \textit{physical states}. The causal continuity is physical, not computational. For machines, there is no such physical continuity; there are only state transitions that are interpreted \textit{as} computations. The paper's symmetry of cases masks this disanalogy.

\section{The Collapse of Organizational Invariance}

Section 6.3 presents organizational invariance---Chalmers's argument from fading and dancing qualia---as the optimist's strongest card \citep{chalmers1996}. But the argument does not reach machine consciousness, and the paper's own qualifications make this clear.

The fading-qualia argument runs: if we replace neurons one at a time with functionally identical units, consciousness cannot fade (without absurdity) or vanish suddenly, so functional organization must preserve consciousness. This argument has force, but only under heavy constraints that the paper acknowledges but does not adequately emphasize.

First, ``functionally identical'' is doing enormous work, and it is unclear whether the intended equivalence holds across machines. For neurons to be ``functionally identical,'' what matters is that they preserve all consciousness-relevant causal powers. But we do not know what the consciousness-relevant causal powers are. If consciousness depends on the specific biochemistry of neuronal transmission---on ion channels, neurotransmitters, neuromodulators, glial support, metabolic coupling---then a silicon system that merely preserves input/output relations or even local spiking patterns will not be functionally identical in the relevant sense.

Second, even granting full functional equivalence, the argument applies only to \textit{replacement within a living system}. The paper makes this explicit: ``a silicon prosthesis embedded in a living brain inherits that brain's organization; a server farm running a language model does not.'' This is the key insight, and it should have been emphasized earlier. Organizational invariance protects only \textit{substitution within a biological context}. Once the biological context is removed, the argument provides no support for consciousness in the resulting system.

Why does the biological context matter? Because consciousness may not be a local property of a single module or even of a single brain. It may be a property of the whole organism---the integration of brain, body, environment, and evolutionary history. If consciousness depends on this larger integration, then we can replace neurons within that system while preserving consciousness. But creating a new system with the same local organization in isolation will not create consciousness, because the larger integrative context is absent.

This is the enactivist insight, and organizational invariance, even if granted, does not refute it. The paper quotes the enactivist: consciousness might be ``dynamical rather than merely computational; embodied rather than intracranial; homeostatic and self-maintaining rather than representational; autopoietic; developmentally and historically formed rather than snapshot-instantiable.'' Organizational invariance says: if you preserve the right local functional organization, consciousness is preserved. Enactivism says: consciousness is not a local property; it depends on the organism's integration with its body and environment over developmental time. These are not refuted by organizational invariance; they are only disputed.

The paper's treatment of enactivism is serious, but the three pressures it identifies do not resolve the debate. Over-generation (all living cells are autopoietic) shows that enactivism needs to specify \textit{which} organizational properties matter, not that the view is wrong. Smuggling the hard problem shows that enactivism faces the same explanatory gap as everyone else, not that it is mistaken. Reduction to other horns shows that enactivism might be wrong, but not that it is wrong. And the paper explicitly does not claim enactivism is wrong. It claims enactivism is an ``unfinished research program rather than a demonstrated impossibility.''

But if consciousness depends on the kind of dynamical, embodied, developmental integration that enactivism describes, then machine consciousness is not merely unproven; it is principally impossible \textit{for systems that lack this integration}. Current digital machines---and foreseeable successors---lack this integration. They are designed systems created ex nihilo, not developed organisms emerging from evolutionary continuity with a body and environment. They are designed for task performance, not for autonomous survival and self-maintenance. They lack the metabolic embedding that characterizes life.

This is not a mystical or obscure claim. It is specific: consciousness may require the kind of causal integration characteristic of \textit{living systems}---systems that maintain themselves through continuous metabolic activity, that develop through sensitive periods in response to environmental interaction, that possess unified physiological regulation, and whose mental capacities coevolved with their embodied needs and capacities. Digital machines lack this integration. Even if we could perfectly replicate neural organization, we would not thereby replicate consciousness, because we would not have replicated the organism.

\section{Biological Substrate and the Irrelevance of ``Base Rates''}

The paper's Section 4 concedes that the biological base rate is ``genuine evidence'' the optimist must answer. It states: ``Every uncontroversial instance of consciousness we possess\ldots arises in biological organisms. That is not a proof of biological essentialism, but it is genuine evidence.''

This concession understates the force of the observation. The appropriate question is not: ``Is this evidence of biological essentialism?'' but rather: ``What should we rationally infer from the fact that \textit{all} conscious systems we encounter are biological, and \textit{no} conscious systems we encounter are non-biological?''

The paper's appeal to Bayes's theorem is correct as far as it goes: a rational agent should weight the observation of an all-biological base rate. But the paper does not adequately discuss what prior probability is warranted for ``machine consciousness'' when we have no positive theory supporting it and no positive case instantiating it. If the prior is very low---say, $0.01$ or lower---and the evidence is entirely from biological systems (100\% biological, 0\% non-biological), then Bayesian updating will not significantly move the posterior probability upward. The observation confirms what we should already expect if machines are not inherently conscious-capable.

More fundamentally, the paper treats biological essentialism as a possible position that might be true or false---a view that needs further evidence to be confirmed or disconfirmed. But it is worth asking whether the burden of proof has been placed correctly. We have overwhelming evidence that biological systems are conscious (or at least, that humans are, and that many animals are). We have no evidence that non-biological systems are conscious. In the absence of a positive argument for why non-biological systems should be conscious-capable, the rational default position is that consciousness is biologically dependent.

This is not mysticism. It is methodologically conservative. We should not posit new consciousness-supporting substrates without empirical evidence or a compelling theoretical argument. The paper attempts to provide the theoretical argument (organizational invariance), but as we have seen, that argument applies only to biological substitution. Until we have either empirical evidence of machine consciousness or a theoretical argument showing why machines should be consciousness-capable, the default is that they are not.

\section{The Identity Problem and Phenomenal Unity}

The paper's Section 4 on identity deserves closer scrutiny. The author invokes Parfit to argue that what matters is psychological continuity, not numerical identity, then claims this leaves identity equally open for all three cases (vat-brain, RL system, us in a discrete cosmos) \citep{parfit1984}.

But this move conflates three distinct questions:

\begin{itemize}
\item \textit{Personal identity}: Is the same person present over time?
\item \textit{Subject identity}: Is there one experiencing subject rather than a succession?
\item \textit{Phenomenal unity}: Is there a unified stream of experience, a ``what it is like'' for the subject?
\end{itemize}

Parfit addresses personal identity and, crucially, argues that the further fact of identity may not matter metaphysically---what matters is psychological continuity. But Parfit explicitly sets aside the question of whether psychological continuity suffices for phenomenal consciousness. And he does not address phenomenal unity directly.

Now, consider an RL system where each step retrains the weights. The paper claims identity is ``equally undetermined'' here as in biological cases, because Parfit helps with the identity-counting puzzle. But Parfit's help is precisely with \textit{personal} identity---with whether to count two psychologically continuous stages as the same person. For phenomenal consciousness, a different question arises: is there \textit{any} unified experience at all, or merely a succession of momentary experiences with no binding relation?

The human case is clear: we have phenomenal unity. Our experiences are unified into a stream of consciousness. This unity may depend on ongoing causal integration, on the persistence of neural structures that maintain phenomenal binding, on the organism's continuity over time. An RL system that is reset between steps loses any such binding; each activation is phenomenologically isolated. The fact that Parfit is agnostic about whether this matters to personal identity does not mean it does not matter to consciousness.

In fact, the paper's invocation of Parfit may have the opposite effect from what is intended. If Parfit is right that personal identity is less metaphysically deep than we assumed, and if consciousness is even more metaphysically opaque than personal identity, then we have reason to be cautious about assuming consciousness in systems that lack the phenomenal unity characteristic of human experience.

\section{The Structure of the Argument and the Burden of Proof}

Stepping back, we can see that the paper's structure shifts the burden of proof in problematic ways. It opens by establishing epistemic possibility---that we lack a deductive proof of impossibility. It then argues that several purported barriers (timing, implementation, computation) do not definitively rule out machines. It acknowledges that biological essentialism and theological positions remain unrefuted. Finally, it defends organizational invariance as the ``better-supported speculation'' against its denying counterpart.

But none of this establishes that we have \textit{positive reason} to expect machine consciousness. It merely establishes that we lack a \textit{knock-down proof} against it. The paper concludes with the ``honest ceiling'': ``nothing we can presently show makes machine consciousness impossible; a leading argument makes it plausible for functional duplicates; the purely computational case\ldots remains open.''

This is a strange place to land. It is true that we cannot prove machine consciousness is impossible. It is also true that we cannot prove it is possible. The paper claims organizational invariance makes consciousness ``plausible for functional duplicates,'' but plausibility is not established by defending an argument against objections; it is established by providing positive reason to expect it. Organizational invariance, even fully granted, does not tell us to expect consciousness in machines; it merely says that \textit{if} a machine were functionally identical to a brain, we might expect consciousness. But we have no independent reason to expect functional identity to obtain across such different substrates, nor do we have reason to expect that functional identity in isolation (divorced from biological embedding) suffices for consciousness.

The deepest issue is this: the paper occupies an unstable middle position. It claims machine consciousness is possible, yet provides no positive theory of consciousness that predicts it. It claims to neutralize objections, yet acknowledges that the most serious objections (biological essentialism, enactivism) survive. It claims to establish modal possibility, yet leaves the question genuinely open. What does it mean for something to be ``possible'' if all our theories and evidence point away from it, all the most serious objections to it are unrefuted, and we have no positive reason to expect it? At some point, ``possible in principle'' becomes indistinguishable from ``merely conceivable'' or ``not demonstrably incoherent''---and conceivability is not a reliable guide to metaphysical possibility.

\section{Consciousness as a Biological Phenomenon}

Let me articulate the positive view that the paper fails to adequately consider. Consciousness, on this view, is not a local computational property. It is a property of \textit{living systems}---systems that maintain themselves through continuous metabolic activity, that develop through sensitive interaction with their environment, that possess unified physiological regulation, and that have coevolved with the embodied needs characteristic of animal life.

This view has several components:

\begin{enumerate}
\item \textbf{Metabolic integration.} Consciousness does not arise in isolated modules; it arises in organisms that regulate their own existence through metabolic coupling with their environment. The brain is not a computer; it is an organ embedded in a body that is embedded in an environment. Its organization reflects millions of years of selection for survival and reproduction in that environment. A digital system, by contrast, requires external maintenance; it cannot sustain itself.

\item \textbf{Developmental history.} Consciousness arises through development. The brain's organization is not static, encoded in a snapshot. It is sculpted through sensitive developmental periods, through embodied interaction with the world, through the gradual myelination and pruning and remodeling that characterizes neural development. A system instantiated ex nihilo, even if perfectly organized, lacks this developmental history. It is not a brain; it is a replica.

\item \textbf{Embodied coupling.} Consciousness is inseparable from embodiment. The brain evolved to control a body in an environment. Its organizational principles---attention, prediction, emotion, motivation---are all structured around the control of an embodied agent. A server farm with no body and no evolved relationship to environmental demands is not the kind of system consciousness could emerge from. It is designed for task performance, not for survival.

\item \textbf{Evolutionary continuity.} Consciousness did not arise ex nihilo in humans. It has precursors in animal life. The phenomenology of human consciousness---the sense of having a body, of being situated in space, of having needs and desires---all reflect our evolutionary continuity with other animals. Digital systems have no evolutionary history. They have no ancestors with simpler forms of consciousness. They represent a radical break from the rest of life.
\end{enumerate}

These claims are not mystical. They are grounded in comparative biology, in developmental neuroscience, in evolutionary theory, in the philosophy of life---Evan Thompson's work, for instance, has shown that autonomy and adaptivity are fundamental to life and to mind \citep{thompson2007}. They suggest that consciousness is not a property that can be instantiated in arbitrary physical systems that happen to have the right formal organization. Consciousness is a biological phenomenon through and through.

The paper takes these possibilities seriously in its discussion of enactivism. But it does not adequately engage with the possibility that they might be right. It treats them as ``research programs'' that remain undecided rather than as powerful and coherent accounts of consciousness that are consistent with all evidence we have and with our best biological understanding. The paper asks: ``What makes consciousness possible?'' and concludes: ``We do not know, so anything is possible.'' But we could instead ask: ``What is consciousness in reality?'' and conclude: ``It is a biological phenomenon,'' with profound implications for its realizability in non-biological substrates.

\section{The Explicitness That Still Conceals}

I close with a methodological critique. The paper is admirably explicit about what it does and does not establish. It separates epistemic from metaphysical possibility. It acknowledges that architectural objections are defeated filters, not proofs of consciousness. It concedes that enactivism and biological essentialism are not refuted. It notes that organizational invariance applies only to functional duplicates, not to arbitrary machines. This explicitness is philosophically valuable.

But the paper's careful qualifications can obscure what is actually being claimed. The central claim is that machine consciousness is ``possible''---that is, metaphysically possible or nomologically possible. Yet the paper's arguments establish, at most, that we lack a deductive proof of impossibility. This is importantly weaker than possibility. Many things are conceivable but not possible. The distinction between failing to rule something out and actively supporting it is crucial, and the paper's language risks obscuring it.

Specifically:

\begin{itemize}
\item The defeater-filter test removes certain objections but does not establish consciousness in machines; it presupposes assumptions about multiple realizability it should defend independently.
\item The vat-brain case shows consciousness survives in biological systems under architectural manipulation; it does not show consciousness would arise in non-biological systems.
\item The discrete-time argument shows consciousness is compatible with discrete temporal structure; it does not show consciousness arises in computational systems that exploit discreteness as a fundamental feature.
\item Organizational invariance shows that \textit{if} a system is a functional duplicate of a brain \textit{and} embedded in the right biological context, consciousness is preserved; it does not show consciousness arises in isolated digital systems.
\end{itemize}

The cumulative effect of these qualifications is that the case for machine consciousness relies entirely on the organizational-invariance argument applied to functional duplicates. But functional duplicates, at full specification, may be indistinguishable from biological brains---in which case they are not really ``machines'' in the sense relevant to the debate. They are biological substitutes. And if consciousness requires biological substrate, then building artificial brains---actual matter structured like biological brains---would not constitute machine consciousness; it would constitute bioengineering.

\section{Conclusion: In Principle Barred}

Machine consciousness is not merely improbable or unproven. It is barred in principle---not by any specific feature machines have (that could be refuted by the defeater-filter test) but by what machines \textit{lack}: the metabolic integration, developmental history, embodied coupling, and evolutionary continuity characteristic of living systems.

The paper's modesty about what can be claimed is admirable. But modesty about what we can prove should not be confused with modesty about what is true. It is true that we cannot deduce consciousness's impossibility in machines from first principles. But it is also true that all our evidence, theory, and understanding suggest consciousness is a biological phenomenon. In the absence of positive evidence to the contrary, the rational stance is not ``possible but unproven,'' but ``impossible given what we know about consciousness and life.''

The paper's final word is that the question ``remains open.'' I propose we should close it: machine consciousness is not open. It is ruled out by principle---not the principle of deductive logic, but the principles of biology and the science of consciousness. Consciousness is what living systems do, and digital machines, no matter how sophisticated, are not living systems. No amount of functional organization can substitute for the integration that characterizes life.

\newpage
\begin{thebibliography}{99}

\bibitem[Anonymous(2026)]{intermittent}
Anonymous. 2026. ``Intermittent Minds: Discreteness, Substrate, and the Possibility of Machine Consciousness.'' Revised manuscript.

\bibitem[Chalmers(1996)]{chalmers1996}
Chalmers, David J. 1996. \textit{The Conscious Mind: In Search of a Fundamental Theory}. Oxford University Press.

\bibitem[Parfit(1984)]{parfit1984}
Parfit, Derek. 1984. \textit{Reasons and Persons}. Oxford University Press.

\bibitem[Searle(1992)]{searle1992}
Searle, John R. 1992. \textit{The Rediscovery of the Mind}. MIT Press.

\bibitem[Thompson(2007)]{thompson2007}
Thompson, Evan. 2007. \textit{Mind in Life: Biology, Phenomenology, and the Sciences of Mind}. Harvard University Press.

\end{thebibliography}

\end{document}